\documentclass[11pt,letterpaper]{article}
\usepackage[T1]{fontenc}
\usepackage[utf8]{inputenc}
\usepackage{lmodern,microtype}
\usepackage[margin=1in,headheight=14pt]{geometry}
\usepackage{amsmath,amssymb,amsthm,mathtools}
\usepackage{booktabs,array,longtable,enumitem}
\usepackage{xcolor}
\usepackage{fancyhdr}
\usepackage[colorlinks=true,linkcolor=blue!45!black,citecolor=green!35!black,urlcolor=blue!45!black]{hyperref}
\usepackage{bookmark}
\hypersetup{pdftitle={Global Hahn Support and the Surcomplex Riemann--Hilbert Problem},pdfauthor={Research manuscript prepared for the Surreal project},pdfsubject={Support-sensitive realization of noncommutative monodromy}}
\pagestyle{fancy}
\fancyhf{}
\fancyhead[L]{\small Global Hahn support and monodromy}
\fancyhead[R]{\small Research manuscript}
\fancyfoot[C]{\thepage}
\setlength{\parskip}{3pt}
\setlength{\parindent}{1.2em}
\setlist{itemsep=3pt,topsep=5pt}
\emergencystretch=2em
\numberwithin{equation}{section}
\newtheorem{theorem}{Theorem}[section]
\newtheorem{lemma}[theorem]{Lemma}
\newtheorem{proposition}[theorem]{Proposition}
\newtheorem{corollary}[theorem]{Corollary}
\theoremstyle{definition}
\newtheorem{definition}[theorem]{Definition}
\newtheorem{example}[theorem]{Example}
\theoremstyle{remark}
\newtheorem{remark}[theorem]{Remark}
\newcommand{\C}{\mathbb C}
\newcommand{\Q}{\mathbb Q}
\newcommand{\Z}{\mathbb Z}
\newcommand{\N}{\mathbb N}
\newcommand{\No}{\mathbf{No}}
\newcommand{\SC}{\mathbf{SC}}
\newcommand{\OO}{\mathcal O}
\newcommand{\HH}{\mathcal H}
\newcommand{\LL}{\mathcal L}
\newcommand{\mm}{\mathfrak m}
\newcommand{\supp}{\operatorname{supp}}
\newcommand{\Hol}{\operatorname{Hol}}
\newcommand{\Log}{\operatorname{Log}}
\newcommand{\Res}{\operatorname{Res}}
\newcommand{\tr}{\operatorname{tr}}
\newcommand{\id}{\operatorname{id}}
\newcommand{\GL}{\operatorname{GL}}
\newcommand{\SL}{\operatorname{SL}}
\newcommand{\Mat}{\operatorname{M}}
\newcommand{\im}{\operatorname{im}}
\newcommand{\val}{\operatorname{v}}
\newcommand{\coeff}[1]{[t^{#1}]}
\newcommand{\cmon}[1]{\langle #1\rangle}
\newcommand{\one}{I_r}
\newcommand{\K}{K_\Gamma}
\newcommand{\Hp}{\HH_{\Gamma,+}}
\newcommand{\commit}{\texttt{39f2be6667ade51bca2b45daa47e289d69c09764}}

\title{\vspace{-1.1cm}\textbf{Global Hahn Support and the\\ Surcomplex Riemann--Hilbert Problem}\\[7pt]
\large Noncommutative monodromy, countably many singularities,\\and an exact realization obstruction}
\author{Research manuscript prepared for the \texttt{Surreal} project}
\date{September 21, 2026}

\begin{document}
\maketitle
\begin{abstract}
Let $D\subset\C$ be closed and discrete, let $U=\C\setminus D$, and let
$K_\Gamma=\C((t^\Gamma))$ be a set-sized Hahn field. We study matrix connections
on $U$ whose coefficients are ordinary holomorphic one-forms, whose Hahn
exponents are positive, and whose support is globally well ordered. For a
representation $\rho:\pi_1(U,p)\to I_r+\Mat_r(\mathfrak m_\Gamma)$, we prove an
exact criterion: a connection of this kind has monodromy $\rho$ if and only if
$\bigcup_{d\in D}\supp(\rho(\ell_d)-I_r)$ is well ordered, where $\ell_d$ are
based meridians. When the criterion holds, every coefficient can be chosen
meromorphic on $\C$ with at most simple poles in $D$, with support in the
additive monoid generated by the prescribed monodromy support. A chosen
linear Mittag--Leffler period section produces a unique normalized realizer.
We also prove framed gauge classification, coefficientwise removability of
gauges between logarithmic realizers, a positive-support Frobenius normal
form, and a determinant-one refinement. An explicit $\SL_2$ representation
has every finite subproblem realizable but admits no global realization;
another example shows that point-finiteness across punctures is not necessary.
A noncommuting two-puncture calculation exhibits the required mixed
commutator correction, including in a genuinely non-Archimedean value group.
These are proposed original results in a specified Hahn-coherent category,
not a claimed resolution of the unrestricted classical Riemann--Hilbert
problem. The main proofs are given in full; the classical inputs and the
limits of the literature and computational checks are stated explicitly.
\end{abstract}

\noindent\textbf{Keywords.} Surcomplex numbers; Hahn series; monodromy;
Riemann--Hilbert problem; Mittag--Leffler theorem; ordered support;
logarithmic connections; non-Archimedean analysis.

\clearpage
\tableofcontents
\clearpage

\section{The problem, the contribution, and its status}\label{sec:intro}

\subsection{Why monodromy is a substantive next question}
A local solution of a complex differential equation may return to a different
branch after continuation around a puncture. In a matrix equation, these
branch changes are matrices and need not commute. The inverse problem asks
whether prescribed branch changes arise from one global differential
equation. This is the monodromy realization, or Riemann--Hilbert, viewpoint.
Classical regular-singular theory provides the background
\cite{deligne,ilyashenko}; our problem is different in its coefficient
category and in the possibility of infinitely many punctures.

The field $\SC=\No[i]$ alone does not specify an analytic category. In the
category used here, a function is a Hahn series whose coefficients are
ordinary holomorphic functions on \emph{one common ordinary domain}.
Analytic continuation is performed coefficientwise. An individual matrix
with infinitesimal entries is legitimate whenever its entries have Hahn
supports. A whole family of such matrices, however, need not have a common
well-ordered support. The inverse-monodromy problem detects this distinction
exactly.

The principal result is not merely a matrix version of a single local
existence theorem. It combines a global condition across countably many
punctures, genuinely noncommuting holonomy, an arbitrary-rank ordered value
group, a constructive transfinite normal form, and a support-preserving gauge
classification. The condition turns out to be both necessary and sufficient.
It also produces an explicit failure of the inference that solvability of
all finite subproblems implies global solvability.

\subsection{The theorem in one display}
Choose a base point $p\in U=\C\setminus D$ and free based meridians
$\ell_d$. If $M_d=\rho(\ell_d)$, put
\begin{equation}\label{eq:headline-support}
 E_\rho=\bigcup_{d\in D}\supp(M_d-I_r).
\end{equation}
The union includes all matrix entries. For a positive well-ordered set $E$,
write $E^*$ for its additive monoid, including $0$. The main theorem proves
\begin{equation}\label{eq:headline}
\begin{gathered}
 E_\rho\text{ is well ordered}\\[-2pt]
 \Updownarrow\\[-2pt]
 \rho=\Hol_\Omega\text{ for a global positive Hahn connection }\Omega,\\
 \text{with }\supp\Omega\subset E_\rho^*\setminus\{0\}
 \text{ and coefficientwise simple poles in }D.
\end{gathered}
\end{equation}
The conclusion is stronger than just existence of some holomorphic
connection on the punctured domain. No additional finite singularities are
introduced. For infinitely many punctures, there is deliberately no
regularity or growth condition at infinity.

\subsection{Relationship to the supplied repository}
The source repository was inspected at commit
\begin{center}\small\commit.\end{center}
Its documentation catalogue and the descriptions of the analysis,
analytic-geometry, global-divisor, and contour packages were reviewed
\cite{repo-catalogue,repo-analysis,repo-geometry,repo-global,repo-contours}.
The existing global-divisor package already studies global well-ordering
obstructions for scalar interpolation and moving divisors. The contour
package already supplies a coefficientwise differential and integration
viewpoint. Those are antecedents, not contributions claimed here.

The present manuscript asks instead for a matrix connection with a prescribed
\emph{nonabelian representation}, proves a necessary-and-sufficient condition,
and constructs a normalized inverse to holonomy. It does not repeat the
repository's Nullstellensatz, finite-zero geometry, residue duality,
trigonometry, or scalar Picard-group calculations. Neither those draft ring
theorems nor a sheaf-cohomology theorem for a Hahn coefficient sheaf is used
in our proofs.

This was not an exhaustive audit of every source manuscript in the
repository. A repository code-search response was marked incomplete, so its
empty result is not evidence of absence. The comparison above is a comparison
with the documented contents actually reviewed.

\subsection{Novelty and the classical ingredients}
The proposed contribution is the precise uniformly supported, countably
punctured, noncommutative realization-and-classification package. The
literature search did not identify this exact theorem, but it does not
establish priority. In particular, monodromy, iterated integrals,
regular-singular equations, Mittag--Leffler interpolation, and recursive
formal deformation methods are classical. Neumann's support calculus is
classical as well \cite{neumann,poonen}; it is not renamed as a new theorem.

M\"uller--Strohmaier study convergent generalized expansions of functions near
zero on a logarithmic cover \cite{muller}. Here $t$ is an independent Hahn
parameter and $z$ is the ordinary complex analytic variable. We do not
identify their analytic category with ours. Deligne's regular-singular
correspondence \cite{deligne} is likewise a major classical reference point,
not a theorem reproved here in its original generality. The familiar
finite-puncture, near-identity formal inversion is the closest classical
shape of the finite special case. The infinite-puncture support obstruction
is the feature on which our novelty assessment chiefly rests.

The manuscript is an AI-assisted research draft. Its general theorems have
not been checked in a proof assistant or refereed. The accompanying program
checks a finite example, not the universal statements or priority. No named
published open conjecture is claimed solved.

\section{The exact coefficient category}\label{sec:category}

\subsection{Hahn workspaces and the surcomplex interpretation}
Fix a set-sized linearly ordered abelian group $\Gamma$. Its Hahn field is
\[
 \K=\C((t^\Gamma))
 =\left\{\sum_{\gamma\in S}c_\gamma t^\gamma:
              S\subset\Gamma\text{ well ordered}\right\}.
\]
The valuation of a nonzero series is its least exponent. Set
\[
 \mm_\Gamma=\{a\in\K:\val(a)>0\}\cup\{0\}.
\]
When $\Gamma$ is an ordered subgroup of $\No$, the normal-form interpretation
$t^\gamma=\omega^{-\gamma}$ embeds the real Hahn workspace in $\No$ and its
complexification in $\SC$. Normal forms and Hahn fields are the background
for this interpretation \cite{poonen,bournez}. Divisibility of $\Gamma$ and
algebraic closedness of $\K$ are not needed below. No construction uses a
proper-class sum, a proper-class fundamental group, or a proper-class ring
of matrix coefficients.

\begin{definition}[Global common-domain Hahn families]\label{def:H}
For an ordinary complex domain $V$, define
\[
 \HH_\Gamma(V)=\OO(V)((t^\Gamma))
\]
to consist of families $F=\sum_\gamma t^\gamma F_\gamma$ with
$F_\gamma\in\OO(V)$ and globally well-ordered support
$\{\gamma:F_\gamma\not\equiv0\}$. Its positive part $\Hp(V)$ consists of
families supported in $\Gamma_{>0}$. The same notation with
$\Omega^1_{\mathrm{hol}}(V)$ in place of $\OO(V)$ denotes one-form families.
For matrices, support is the finite union of the entry supports.
\end{definition}

A positive connection on the trivial rank-$r$ bundle over $U$ means
\begin{equation}\label{eq:connection}
 dY=\Omega Y,\qquad
 \Omega=\sum_{\gamma\in S}t^\gamma\Omega_\gamma,
 \qquad S\subset\Gamma_{>0}\text{ well ordered},
\end{equation}
where each $\Omega_\gamma$ is an ordinary holomorphic matrix one-form on
\emph{all} of $U$. In complex dimension one these holomorphic connections
are flat: both the exterior derivative of a holomorphic one-form and the
wedge of two such one-forms vanish.

The word ``positive'' always concerns exponents, not positive-definite
matrices. Similarly, $I_r+\Mat_r(\mm_\Gamma)$ is a group of matrices
congruent to the identity. Its elements need not be algebraically unipotent.

\subsection{What common domain and global support exclude}
The objects in Definition~\ref{def:H} are not arbitrary formal series in $z$,
not arbitrary families of individually convergent germs, and not arbitrary
local sections with unrelated supports. They are also generally not the
algebraic tensor product $\K\otimes_\C\OO(V)$: a Hahn family may involve
infinitely many linearly independent holomorphic coefficient functions.

We do not assume that the global-support assignment on all open sets is a
sheaf for arbitrary covers. All our gluing is proved directly by descending
families already supported in one specified monoid. The only use of
Mittag--Leffler is an ordinary complex theorem, one exponent at a time.
Consequently no application of Cartan's theorem to an infinite-dimensional
Hahn coefficient sheaf is hidden in the argument.

\subsection{Ordinary paths, not fine-surreal paths}
For a piecewise smooth ordinary path $c$ in $U$, integrals of
$\Omega_\gamma$ and their iterated integrals are ordinary complex integrals.
Their values are then assembled as Hahn coefficients. We do not assert
convergence of partial Hahn sums in the fine topology of $\SC$. This is the
same categorical separation emphasized by the supplied analysis and contour
reports \cite{repo-analysis,repo-contours}.

For completeness, the coefficient category has its usual infinitesimal
surcomplex interpretation. If $z_0\in V$ and
$\varepsilon\in\mm_\Gamma$, then
\begin{equation}\label{eq:halo-eval}
 F^\sharp(z_0+\varepsilon)=
 \sum_{\gamma,k\geq0}
 t^\gamma\frac{F_\gamma^{(k)}(z_0)}{k!}\varepsilon^k
\end{equation}
for a positive family, with the ordinary exponent-zero term allowed
separately. The sum is strongly summable by the support lemma below, applied
to the union of the positive coefficient support and the support of
$\varepsilon$. Each coefficient in a fixed Hahn exponent receives only
finitely many contributions. This explains why the construction belongs to
surcomplex analysis, but the proof of monodromy is carried out on the
ordinary base $U$ and its ordinary universal cover. Monodromy is not defined
by transporting along a fine-continuous surcomplex loop.

\section{Support calculus without a convergence assumption}\label{sec:support}

\begin{definition}[Strong summability]
A set-indexed family of Hahn series is strongly summable if the union of its
supports is well ordered and each exponent occurs in the support of only
finitely many members. For vector- or matrix-valued coefficients, the same
definition is used. The sum is then the finite coefficientwise sum at each
exponent.
\end{definition}

\begin{lemma}[Classical Neumann support lemma]\label{lem:neumann}
Let $E\subset\Gamma_{>0}$ be well ordered. The set
\[
 E^*=\{0\}\cup
 \{e_1+\cdots+e_n:n\geq1,\ e_j\in E\}
\]
is well ordered. For each $\gamma\in\Gamma$, there are only finitely many
ordered finite words $(e_1,\ldots,e_n)$ in $E$ with sum $\gamma$.
Also, the sum of two well-ordered subsets of $\Gamma$ is well ordered and a
fixed element has only finitely many decompositions with one summand in each
of the two sets.
\end{lemma}

We use this standard support theorem as a classical input
\cite{neumann,poonen}. It is stronger than a statement about valuations
increasing along one particular sequence. In particular, the finite-word
assertion applies across all word lengths, which is the point needed for
noncommutative Dyson expansions and geometric inverses.

\begin{lemma}[Algebraic operations with a common support certificate]
\label{lem:operations}
Suppose every entry of a matrix $A$ is supported in a positive well-ordered
set $E$. Then
\[
 (I_r+A)^{-1}=\sum_{n\geq0}(-A)^n,
 \qquad
 \exp A=\sum_{n\geq0}\frac{A^n}{n!},
 \qquad
 \log(I_r+A)=\sum_{n\geq1}\frac{(-1)^{n+1}A^n}{n}
\]
are strongly summable, with support in $E^*$. The same holds for coefficients
in an associative complex algebra, provided the coefficient operations
appearing in each finite expression are defined.
\end{lemma}
\begin{proof}
Expand each matrix entry into products indexed by words in $E$, with the
finitely many matrix indices inserted. A contribution to exponent $\gamma$
can come from only finitely many exponent words by
Lemma~\ref{lem:neumann}; for each word there are only finitely many matrix
index choices. All supports lie in $E^*$. This proves strong summability.
The usual finite geometric identity, examined coefficientwise, gives the
inverse formula because for a fixed exponent sufficiently long words do
not contribute. The formal exponential and logarithm identities used below
are justified in exactly the same way, as coefficientwise finite identities.
\end{proof}

\begin{lemma}[Finite decomposition inside the generated monoid]
\label{lem:finite-closure}
Put $S=E^*\setminus\{0\}$. Each $\gamma\in\Gamma$ has only finitely many
ordered decompositions into members of $S$. Moreover, only finitely many
members of $S$ can occur as a part in any such decomposition of a fixed
$\gamma$.
\end{lemma}
\begin{proof}
Choose a representation by an $E$-word for each member of $S$. Flattening a
word in $S$ of total weight $\gamma$ gives an $E$-word of that weight.
There are finitely many such flattened words. A fixed finite word has only
finitely many partitions into nonempty consecutive blocks. Every original
word in $S$ is recovered from one of these finite possibilities. The second
assertion follows by taking all parts of those finitely many words.
\end{proof}

\begin{remark}[Why a valuation-limit argument would be wrong]
Take $\Gamma=\Q+\Q\omega$ with its induced surreal order. Then $n<\omega$
for every ordinary positive integer $n$. Thus repeatedly increasing a
valuation by $1$ never reaches $\omega$. Nevertheless
$\sum_{n\geq0}t^n$ is a valid Hahn sum, and the support calculus works equally
well at exponents $\omega$ or $\omega+1$. Our recursions are well-founded
recursions on well-ordered support sets, not proofs that repeated positive
increments become cofinal in the value group.
\end{remark}

\section{Ordinary topology and a period section}\label{sec:period}

\subsection{A compatible meridian basis}
Let $D\subset\C$ be closed and discrete. Every compact set meets $D$ in a
finite set, so $D$ is at most countable. The domain $U=\C\setminus D$ is
connected. Choose a base point $p\in U$.

\begin{lemma}[Free based meridians]\label{lem:free}
There is a collection of based positively oriented meridians
$(\ell_d)_{d\in D}$ forming a free basis of $\pi_1(U,p)$. Every group element
is a finite word in these meridians and their inverses.
\end{lemma}
\begin{proof}
Exhaust $\C$ by nested disks centered at $p$, with boundaries missing $D$.
Each disk contains finitely many punctures. A finitely punctured disk has
free fundamental group generated by based meridians. When the disk is
enlarged, the old meridians can be retained and meridians for the new
punctures added: cutting the annular region along arcs to the newly included
punctures, or applying van Kampen, exhibits the old group as a free factor.
Choose these bases compatibly. A loop and any specified homotopy have compact
image and therefore occur in a sufficiently large disk. It follows that the
fundamental group of the union is the free group on the union of the
compatible bases. If $D$ is empty this group is trivial.
\end{proof}

We fix such a basis throughout. The numerical matrices assigned to based
meridians depend on the chosen tails and the base point; that dependence is
not suppressed. To make $dY=\Omega Y$ give a homomorphism with matrix
multiplication, we use the convention that $\alpha\beta$ means traversal of
$\beta$ first and then $\alpha$. All meridians themselves are positively
oriented.

\subsection{The ordinary Mittag--Leffler step}
Let $\LL(D)$ be the complex vector space of meromorphic one-forms on $\C$
with at most simple poles at the points of $D$ and no other finite poles.
There is no condition at infinity. Define
\begin{equation}\label{eq:period-map}
 P:\LL(D)\longrightarrow\C^D,
 \qquad (P\alpha)_d=\int_{\ell_d}\alpha=2\pi i\Res_d\alpha.
\end{equation}
The target is the full product of complex vector spaces, not the direct sum.

\begin{lemma}[A linear period right inverse]\label{lem:period-section}
The map $P$ is surjective. There exists a complex-linear map
$\sigma:\C^D\to\LL(D)$ with $P\sigma=\id$. If $D$ is finite, one may take
\begin{equation}\label{eq:finite-section}
 \sigma((v_d)_{d\in D})=
 \frac1{2\pi i}\sum_{d\in D}v_d\frac{dz}{z-d}.
\end{equation}
For $D=\varnothing$ the right inverse has zero image.
\end{lemma}
\begin{proof}
Prescribe the principal part
$v_d\,dz/(2\pi i(z-d))$ at each $d$. The ordinary Mittag--Leffler theorem
supplies a meromorphic one-form with those principal parts, proving
surjectivity. Here is the elementary mechanism relevant to this use. If
$D$ is infinite, enumerate its points so that they escape each compact set.
For all but finitely many points, subtract a sufficiently long Taylor
polynomial from $v_d/(2\pi i(z-d))$ so that the corrected term is smaller
than $2^{-n}$ on a chosen expanding disk that still avoids $d$. The corrected
series converges normally on compact subsets away from $D$, and its local
principal parts are unchanged. Initial exceptional terms are kept
separately. This is an ordinary complex convergence argument.

The construction just described proves surjectivity but need not be linear
in the input because the chosen polynomial degrees may depend on it. To
obtain a linear section, split the surjective map of complex vector spaces:
choose a basis of $\C^D$, choose one preimage for each basis element, and
extend linearly. This uses the usual axiom of choice for vector spaces.
For finite $D$, the displayed finite sum is an explicit linear section.
\end{proof}

\begin{remark}[What is and is not required of the section]\label{rem:section}
For infinite $D$, we assert neither continuity nor a norm estimate nor an
effective algorithm for $\sigma$. In particular, the formula
$\sum_d v_d/(z-d)$ cannot simply be written without a convergence correction.
The theorem needs only an ordinary complex-linear splitting. The section is
chosen once, independently of the Hahn exponent, and is applied entrywise to
matrix-valued data. This last point is important for the traceless
refinement in Section~\ref{sec:sl}.
\end{remark}

\section{Fundamental solutions and support-controlled holonomy}
\label{sec:holonomy}

\subsection{A global solution on the ordinary universal cover}
Let $\widetilde U\to U$ be the ordinary universal covering, with a chosen lift
$\widetilde p$ of $p$. Pullbacks of coefficients to the cover will be
understood.

\begin{theorem}[Fundamental matrix with a support certificate]
\label{thm:fundamental}
For a connection \eqref{eq:connection}, there is a unique positive
fundamental matrix
\[
 Y=I_r+\sum_{\gamma>0}t^\gamma Y_\gamma
 \quad\text{on }\widetilde U,
 \qquad dY=\Omega Y,\quad Y(\widetilde p)=I_r.
\]
Every coefficient is holomorphic on all of $\widetilde U$, and
\[
 \supp(Y-I_r)\cup\supp(Y^{-1}-I_r)\subset S^*\setminus\{0\}.
\]
Analytic continuation determines a representation
$\Hol_\Omega:\pi_1(U,p)\to I_r+\Mat_r(\mm_\Gamma)$ with the same uniform
support bound for every group element.
\end{theorem}
\begin{proof}
Write $T=S^*\setminus\{0\}$. Define the coefficients by well-founded
recursion on $T$. Missing coefficients of $\Omega$ are zero. At weight
$\gamma$, set
\begin{equation}\label{eq:solution-recursion}
 dY_\gamma=\Omega_\gamma+
   \sum_{\substack{\alpha+\beta=\gamma\\\alpha,\beta>0}}
       \Omega_\alpha Y_\beta,
 \qquad Y_\gamma(\widetilde p)=0.
\end{equation}
Only finitely many pairs contribute, by Lemma~\ref{lem:neumann}, and
$\beta<\gamma$. The right-hand side is a holomorphic matrix one-form on the
simply connected Riemann surface $\widetilde U$. Its entries have unique
holomorphic primitives vanishing at $\widetilde p$. This constructs
$Y_\gamma$ on the whole cover. The resulting Hahn family has support in the
well-ordered set $T$ and satisfies the equation coefficientwise.
Lemma~\ref{lem:operations} gives its inverse with the same monoid support.

To prove uniqueness even against a solution with a different positive
well-ordered support, take the least exponent at which the two solutions
differ, in the well-ordered union of their supports. At this exponent their
difference has derivative zero and value zero at $\widetilde p$; it is zero,
a contradiction.

A continued fundamental matrix differs from the original branch on a
connected overlap by a constant matrix on the right. Indeed,
$d(Y^{-1}\widehat Y)=0$ follows coefficientwise from the differential
equation. Continuing around a based loop gives the normalized monodromy
matrix, whose coefficients are values of the coefficients of $Y$. Therefore
all its positive exponents belong to $T$, independently of the loop.
Uniqueness of continuation gives the group law with the convention fixed
in Section~\ref{sec:period}. Homotopy invariance follows from the universal
cover construction.
\end{proof}

\subsection{The explicit iterated-integral expression}
For a piecewise smooth path $c:[0,1]\to U$, write
$c^*\Omega_\gamma=A_\gamma(s)\,ds$. Its transport is
\begin{equation}\label{eq:dyson}
 T_\Omega(c)=I_r+
 \sum_{n\geq1}\ \sum_{\gamma_1,\ldots,\gamma_n\in S}
 t^{\gamma_1+\cdots+\gamma_n}
 \int_{0<s_1<\cdots<s_n<1}
 A_{\gamma_n}(s_n)\cdots A_{\gamma_1}(s_1)\,ds_1\cdots ds_n.
\end{equation}
This is the usual ordered-integral solution, a classical mechanism related
to iterated path integrals \cite{chen}. Here its meaning is particularly
strict: the outer family is strongly summable by the finite-word part of
Lemma~\ref{lem:neumann}. Each individual coefficient is a \emph{finite} sum
of ordinary iterated integrals. The expression agrees with
Theorem~\ref{thm:fundamental}, because termwise differentiation gives
\eqref{eq:solution-recursion}. For a loop it gives $\Hol_\Omega$.

\begin{lemma}[Triangular change of monodromy]\label{lem:triangular}
Suppose two positive connections $\Omega$ and $\Psi$ agree at every exponent
strictly below $\gamma$. Then their transports along any fixed path agree
below $\gamma$, and for a based loop $\ell$,
\begin{equation}\label{eq:triangular}
 \coeff{\gamma}\bigl(\Hol_\Psi(\ell)-\Hol_\Omega(\ell)\bigr)
       =\int_\ell(\Psi_\gamma-\Omega_\gamma).
\end{equation}
Terms strictly above $\gamma$ cannot affect coefficients at or below
$\gamma$.
\end{lemma}
\begin{proof}
Compare the word expansions \eqref{eq:dyson}, using the well-ordered union of
the two positive supports. A differing coefficient of weight at least
$\gamma$ cannot occur in a word of total weight less than $\gamma$. At
weight exactly $\gamma$, a differing coefficient can occur only as the
single letter of the word: every additional letter has strictly positive
weight. The length-one term is the displayed ordinary period. All
comparisons are finite at each exponent.
\end{proof}

The crucial feature is that the linearized map at \emph{every} successive
weight is the same ordinary period map $P$. Noncommutativity appears in the
already determined lower-weight terms; it does not change the operator
needed to solve the new correction equation.

\section{The exact realization theorem}\label{sec:main}

\begin{definition}[Uniform support admissibility]
A representation
$\rho:\pi_1(U,p)\to I_r+\Mat_r(\mm_\Gamma)$ is uniformly support-admissible
if the set $E_\rho$ in \eqref{eq:headline-support} is well ordered.
For a fixed period section $\sigma$, a connection is
\emph{$\sigma$-normalized} if every entry of every coefficient belongs to
$\im\sigma\subset\LL(D)$.
\end{definition}

\begin{theorem}[Support-sensitive Riemann--Hilbert realization]
\label{thm:RH}
Let $D\subset\C$ be closed and discrete, let $p\in U=\C\setminus D$, and fix
free based meridians. For
$\rho:\pi_1(U,p)\to I_r+\Mat_r(\mm_\Gamma)$, the following are equivalent:
\begin{enumerate}[label=\textnormal{(\roman*)}]
 \item $E_\rho$ is well ordered;
 \item $\rho$ is the monodromy of a positive globally common-domain
 Hahn connection on $U$;
 \item $\rho$ is the monodromy of such a connection whose coefficients are
 meromorphic on $\C$ with at most simple poles in $D$ and no other finite
 poles.
\end{enumerate}
If these conditions hold, then for every fixed complex-linear period section
$\sigma$ there is exactly one $\sigma$-normalized realizer $\Omega^\sigma$.
It satisfies
\begin{equation}\label{eq:RH-support}
 \supp\Omega^\sigma\subset E_\rho^*\setminus\{0\}.
\end{equation}
Uniqueness is among all positive globally supported
$\sigma$-normalized connections, not just those whose support was assumed
in advance to lie in the right-hand side of \eqref{eq:RH-support}.
\end{theorem}

\begin{proof}
\emph{Necessity.} Suppose $\rho=\Hol_\Omega$ and $\supp\Omega=S$.
Theorem~\ref{thm:fundamental} gives
$\supp(\rho(g)-I_r)\subset S^*\setminus\{0\}$ for every group element $g$.
Thus $E_\rho$ is a subset of one well-ordered set. This proves
(ii)$\Rightarrow$(i), and (iii)$\Rightarrow$(ii) is immediate.

\emph{Construction.} Assume (i), put
$T=E_\rho^*\setminus\{0\}$, and write
\[
 M_d=\rho(\ell_d)=I_r+\sum_{\gamma\in T}t^\gamma m_{d,\gamma}.
\]
Missing coefficients are zero. We recursively construct
$\Omega_\gamma\in\Mat_r(\LL(D))$ for $\gamma\in T$. Suppose all lower
coefficients have been constructed and form the legitimate positive Hahn
connection
\[
 \Omega^{<\gamma}=\sum_{\substack{\delta\in T\\\delta<\gamma}}
                            t^\delta\Omega_\delta.
\]
Its support is a well-ordered subset of $T$; the initial segment need not be
finite. Define the lower-order prediction and the remaining defect by
\begin{equation}\label{eq:defect}
 h_{d,\gamma}^{<\gamma}
       =\coeff{\gamma}\Hol_{\Omega^{<\gamma}}(\ell_d),
 \qquad
 b_{d,\gamma}=m_{d,\gamma}-h_{d,\gamma}^{<\gamma}.
\end{equation}
The prediction exists by Theorem~\ref{thm:fundamental}. Equivalently it is
the finite sum of all iterated-integral words of total weight $\gamma$
formed from previously constructed coefficients. Apply the period section
entrywise and set
\begin{equation}\label{eq:RH-recursion}
 \boxed{\quad \Omega_\gamma
       =\sigma\bigl((b_{d,\gamma})_{d\in D}\bigr).\quad}
\end{equation}
The datum at one exponent is an arbitrary matrix-valued element of the
\emph{product} $\C^D$; Lemma~\ref{lem:period-section} applies even when
infinitely many of its entries are nonzero. This is why no summability
across punctures is required.

The coefficient $\Omega_\gamma$ is meromorphic on $\C$ with at most simple
poles in $D$. It corrects exactly the defect at weight $\gamma$, because
\[
 \int_{\ell_d}\Omega_\gamma=b_{d,\gamma}
\]
and Lemma~\ref{lem:triangular} applies. It leaves all lower weights
unchanged. Transfinite recursion on the well-ordered set $T$ therefore
constructs
\[
 \Omega^\sigma=\sum_{\gamma\in T}t^\gamma\Omega_\gamma.
\]
All coefficients are holomorphic on the same domain $U$, and its support is
a subset of $T$. No limiting assertion at a limit ordinal is involved:
there is simply a coefficient assigned to every member of a well-ordered
set.

To verify the final monodromy, fix $\gamma\in T$. Coefficients at weights
above $\gamma$ cannot affect it by Lemma~\ref{lem:triangular}; the recursion
has already set it equal to $m_{d,\gamma}$. The monodromy support of the
final connection lies in $T^*\setminus\{0\}=T$, so there are no additional
exponents outside $T$. Thus every meridian has its prescribed monodromy.
The meridians freely generate the fundamental group, proving
$\Hol_{\Omega^\sigma}=\rho$ and (iii).

\emph{Uniqueness.} Suppose $\Omega$ and $\Psi$ are two positive
$\sigma$-normalized realizers. If they differ, the union of their supports
is well ordered and there is a least differing exponent $\gamma$.
Their holonomies are equal, so Lemma~\ref{lem:triangular} implies
\[
 P(\Psi_\gamma-\Omega_\gamma)=0.
\]
But $\Psi_\gamma-\Omega_\gamma\in\im\sigma$ entrywise. The identity
$P\sigma=\id$ says that $P$ is injective on this image. Consequently the
difference is zero, a contradiction.
\end{proof}

\begin{remark}[No hidden convergence over the puncture set]
The construction is not a sum of one independently built differential
equation for every puncture. At a fixed Hahn exponent it first combines
\emph{all} puncture data into one ordinary meromorphic one-form using the
period section. Only these ordinary forms are subsequently assembled in
Hahn exponents. This distinction is indispensable when infinitely many
punctures use the same exponent; see Example~\ref{ex:same-exponent}.
\end{remark}

\subsection{An intrinsic version of the support condition}

\begin{proposition}[Independence of the meridian basis]\label{prop:intrinsic}
For a near-identity representation $\rho$, define
\[
 E_{\mathrm{all}}(\rho)=
 \bigcup_{g\in\pi_1(U,p)}\supp(\rho(g)-I_r).
\]
Then $E_\rho$ is well ordered if and only if
$E_{\mathrm{all}}(\rho)$ is well ordered. Thus admissibility does not depend
on the chosen free meridians or, after a base-point identification, the base
point.
\end{proposition}
\begin{proof}
One implication follows from $E_\rho\subset E_{\mathrm{all}}(\rho)$.
For the other, assume $E_\rho$ is well ordered. Each group element is a
finite word in the $M_d$ and their inverses. Lemma~\ref{lem:operations}
expands the inverses in the common monoid $E_\rho^*$, and multiplication
preserves that monoid. Hence
$E_{\mathrm{all}}(\rho)\subset E_\rho^*\setminus\{0\}$.
Changing based meridians or identifying fundamental groups at different
base points changes the presentation of the same representation, not this
condition on all its values.
\end{proof}

\begin{proposition}[Fixed conjugation and enlargement cannot repair support]
\label{prop:conjugation}
For a family of matrices $(M_d)$ over $\K$, the property that
$\bigcup_d\supp(M_d-I_r)$ is well ordered is invariant under conjugation by
any one fixed $C\in\GL_r(\K)$. It is also unchanged when the family is
embedded by an order-preserving value-group embedding into a larger Hahn
workspace. An arbitrary conjugation need not preserve near-identity
positivity; the assertion about well-ordered supports does not require
that additional property.
\end{proposition}
\begin{proof}
Let $A$ and $B$ be the supports of the finitely many entries of $C$ and
$C^{-1}$. They are well ordered. If $E$ is the union of the supports of
$M_d-I_r$, then
\[
 \bigcup_d\supp\bigl(C(M_d-I_r)C^{-1}\bigr)\subset A+E+B.
\]
When $E$ is well ordered, so is the set on the right by the last part of
Lemma~\ref{lem:neumann}. The converse follows by conjugating with
$C^{-1}$. An order embedding preserves and reflects the ordering of a
subset, hence preserves and reflects its being well ordered.
\end{proof}

\section{Framed gauge classification and removable gauges}
\label{sec:gauge}

\begin{definition}[Positive framed gauge]
A positive gauge on $U$ is a matrix
\[
 G\in I_r+\Mat_r(\Hp(U)).
\]
It is framed at $p$ when $G(p)=I_r$ as an equality of Hahn matrices.
Its inverse is a positive globally supported holomorphic matrix by
Lemma~\ref{lem:operations}. It transforms a connection according to
\begin{equation}\label{eq:gauge}
 \Psi=dG\,G^{-1}+G\Omega G^{-1}.
\end{equation}
\end{definition}

\begin{theorem}[Framed classification with controlled support]
\label{thm:gauge}
Two positive globally Hahn-supported connections $\Omega$ and $\Psi$ on $U$
have the same based monodromy if and only if there is a positive framed
gauge $G$ satisfying \eqref{eq:gauge}. Such a gauge is unique. If
$E=\supp\Omega\cup\supp\Psi$, then
\begin{equation}\label{eq:gauge-support}
 \supp(G-I_r)\cup\supp(G^{-1}-I_r)\subset E^*\setminus\{0\}.
\end{equation}
Consequently, holonomy gives a bijection between framed positive gauge
classes and uniformly support-admissible near-identity representations.
\end{theorem}
\begin{proof}
Let $Y_\Omega$ and $Y_\Psi$ be the normalized fundamental matrices on
$\widetilde U$. If their based monodromies agree, form
\[
 \widetilde G=Y_\Psi Y_\Omega^{-1}.
\]
Upon continuation around any loop both fundamental matrices acquire the
same right factor $M$. Therefore the product becomes
$Y_\Psi M(Y_\Omega M)^{-1}=Y_\Psi Y_\Omega^{-1}$.
Each ordinary coefficient descends to $U$, and the whole family already
has support in the monoid $E^*$ by Theorem~\ref{thm:fundamental} and
Lemma~\ref{lem:operations}. Thus the descended matrix is a globally
supported positive gauge $G$, with $G(p)=I_r$. Direct differentiation gives
\[
 dG=\Psi G-G\Omega,
\]
which is equivalent to \eqref{eq:gauge}.

Conversely, if a framed gauge exists, $GY_\Omega$ solves the $\Psi$ equation
and has initial value $I_r$. Uniqueness in
Theorem~\ref{thm:fundamental} gives $GY_\Omega=Y_\Psi$. Since $G$ is
single-valued on $U$, the monodromy matrices agree. This equality also
forces $G=Y_\Psi Y_\Omega^{-1}$, proving uniqueness and the support bound.
The bijection now follows from Theorem~\ref{thm:RH}.
\end{proof}

The logarithmic subclass has a stronger rigidity property: although the
comparison gauge is initially defined only on the punctured plane, no
coefficient of it actually needs a puncture.

\begin{theorem}[Entire extension of a logarithmic comparison gauge]
\label{thm:entire-gauge}
Suppose $\Omega$ and $\Psi$ are as in Theorem~\ref{thm:gauge}, have the same
based monodromy, and every coefficient of both lies in $\Mat_r(\LL(D))$.
Then the unique framed gauge and its inverse extend coefficientwise
holomorphically to all of $\C$. In particular,
\[
 G\in I_r+\Mat_r(\Hp(\C)),
 \qquad G^{-1}\in I_r+\Mat_r(\Hp(\C)),
\]
with the same support bound \eqref{eq:gauge-support}.
\end{theorem}
\begin{proof}
Write $G=I_r+\sum_{\gamma>0}t^\gamma G_\gamma$. The gauge equation at
weight $\gamma$ reads
\begin{equation}\label{eq:gauge-recursion}
 dG_\gamma=\Psi_\gamma-\Omega_\gamma+
 \sum_{\substack{\alpha+\beta=\gamma\\\alpha,\beta>0}}
       (\Psi_\alpha G_\beta-G_\beta\Omega_\alpha).
\end{equation}
Induct on the positive support monoid in \eqref{eq:gauge-support}. If every
lower $G_\beta$ is entire, the right-hand side is an ordinary meromorphic
one-form on $\C$ with at most simple poles in $D$; the sum is finite.
Fix $d\in D$ and a disk containing no other point of $D$. The coefficient
$G_\gamma$ is already single-valued and holomorphic on the punctured disk.
Thus the integral of its derivative around a small circle is zero.
The right-hand side of \eqref{eq:gauge-recursion} has residue zero at $d$.
Since it had at most a simple pole, it is holomorphic there. Its primitive
on the full disk differs from $G_\gamma$ by a constant on the punctured
disk, so $G_\gamma$ extends holomorphically across $d$.

Doing this for all $d$ extends $G_\gamma$ to an entire function. This proves
the induction without any uniform bound on the sizes of the coefficient
functions. The ordinary domain after extension is the same $\C$ for all
coefficients. Finally the geometric inverse of the positive entire family
has entire coefficients and the same monoid support.
\end{proof}

\begin{corollary}[A global normal slice for logarithmic connections]
For a fixed $\sigma$, every positive logarithmic connection has exactly one
$\sigma$-normalized representative under framed positive entire gauges.
The representative is the connection constructed from its monodromy by
\eqref{eq:RH-recursion}.
\end{corollary}
\begin{proof}
Its monodromy is admissible by Theorem~\ref{thm:fundamental}.
Theorems~\ref{thm:RH} and \ref{thm:entire-gauge} give existence; uniqueness
of the normalized connection and of the framed gauge has already been
proved.
\end{proof}

\section{Positive Frobenius form at a puncture}\label{sec:frobenius}

The next theorem explains why residues still control \emph{local} monodromy,
but cannot in general be read directly from the prescribed \emph{based}
monodromy matrices. It is a support-controlled version of the elementary
nonresonant regular-singular construction, not a new claim about classical
Frobenius theory.

\begin{theorem}[A common-disk positive Frobenius factorization]
\label{thm:frobenius}
Let $\Delta$ be an ordinary disk about $0$. Suppose a positive Hahn
connection on $\Delta\setminus\{0\}$ has the form
\begin{equation}\label{eq:local-connection}
 \Omega=R\frac{dz}{z}+B(z)\,dz,
\end{equation}
where $R\in\Mat_r(\mm_\Gamma)$ and
$B\in\Mat_r(\Hp(\Delta))$, with one positive well-ordered support $E$ for
the pair. Then there is a unique
$H\in I_r+\Mat_r(\Hp(\Delta))$, $H(0)=I_r$, such that on a logarithm branch
\begin{equation}\label{eq:frobenius}
 Y(z)=H(z)\exp(R\Log z)
\end{equation}
solves $dY=\Omega Y$. The supports of $H-I_r$ and $H^{-1}-I_r$ lie in
$E^*\setminus\{0\}$. Local monodromy based at $z_0\in\Delta\setminus\{0\}$
is
\begin{equation}\label{eq:local-monodromy}
 H(z_0)\exp(2\pi iR)H(z_0)^{-1}.
\end{equation}
\end{theorem}
\begin{proof}
The required equation for $H$ is
\begin{equation}\label{eq:H-equation}
 H'=\frac{[R,H]}{z}+BH,\qquad H(0)=I_r.
\end{equation}
Construct its coefficients recursively in $E^*\setminus\{0\}$. At weight
$\gamma$ one must solve
\[
 H_\gamma'=B_\gamma+
 \sum_{\substack{\alpha+\beta=\gamma\\\alpha,\beta>0}}
 \left(B_\alpha H_\beta+\frac{[R_\alpha,H_\beta]}{z}\right),
 \qquad H_\gamma(0)=0.
\]
This is a finite sum. The previously constructed $H_\beta$ vanish at $0$,
so $H_\beta/z$ is holomorphic on the entire disk. The right-hand side is
therefore holomorphic on the same disk, and integration from $0$ defines a
unique holomorphic $H_\gamma$ vanishing there. The resulting family has the
claimed support; its inverse is given by
Lemma~\ref{lem:operations}. The least-difference argument proves uniqueness.

Because $R$ has positive support, $\exp(R\Log z)$ is a strongly summable
family of ordinary holomorphic functions on any logarithm branch.
Differentiating it coefficientwise gives
$(\exp(R\Log z))'=R\exp(R\Log z)/z$, and
\eqref{eq:H-equation} verifies \eqref{eq:frobenius}. Continuing once
positively around $0$ increases $\Log z$ by $2\pi i$ and multiplies the
fundamental matrix on the right by $\exp(2\pi iR)$. Normalizing at $z_0$
conjugates this factor by
$H(z_0)\exp(R\Log z_0)$; the exponential factor commutes with
$\exp(2\pi iR)$ and cancels in the conjugation. This gives
\eqref{eq:local-monodromy}.
\end{proof}

There is no division by a spectral resonance operator in this proof.
The coefficient at exponent zero of the residue is zero, so the equation at
each new exponent is an ordinary primitive problem. Positivity is doing
real work here.

If a meridian is based at a distant point $p$, its tail contributes a further
transport conjugation to \eqref{eq:local-monodromy}. With several punctures
these conjugations interact. Thus choosing
$R_d=(2\pi i)^{-1}\log M_d$ separately at each puncture need not solve the
based inverse problem. The explicit failure first appears at a mixed
commutator term in Section~\ref{sec:example}.

\section{Finite-puncture inversion and determinant one}\label{sec:finite}

\subsection{A finite Fuchsian specialization}
\begin{corollary}[Unique positive residue matrices for finite data]
\label{cor:finite}
Let $D=\{a_1,\ldots,a_m\}$ be finite. Fix based meridians and arbitrary
matrices $M_j\in I_r+\Mat_r(\mm_\Gamma)$. There are unique residue matrices
$R_j\in\Mat_r(\mm_\Gamma)$ for which
\begin{equation}\label{eq:fuchsian}
 \Omega=\sum_{j=1}^m R_j\frac{dz}{z-a_j}
\end{equation}
has based meridian monodromies $M_j$. If
$E=\bigcup_j\supp(M_j-I_r)$, then $\supp R_j\subset E^*\setminus\{0\}$.
The form has at most a simple pole at infinity, with residue
$-\sum_jR_j$ there.
\end{corollary}
\begin{proof}
A finite union of well-ordered sets in a linear order is well ordered, so
$E$ is admissible. Use the explicit section \eqref{eq:finite-section} in
Theorem~\ref{thm:RH}; being normalized for this section is precisely the
form \eqref{eq:fuchsian}. Its coefficientwise uniqueness gives uniqueness
of all the $R_j$. In the coordinate $w=1/z$,
$dz/(z-a_j)=-dw/(w(1-a_jw))$, proving the assertion at infinity.
\end{proof}

This is a near-identity Hahn version of a classical Fuchsian inverse problem,
not a realization result for arbitrary complex monodromy on a globally
trivial bundle. The residue matrices must have positive support, and a
regular singularity at infinity is permitted. No assertion is made that the
connection is holomorphic at infinity, nor that arbitrary residue matrices
with the same monodromy are unique. Ordinary integer shifts of exponents,
which already destroy such unrestricted uniqueness, are excluded by the
positive-support hypothesis.

\subsection{The special-linear refinement}\label{sec:sl}

\begin{lemma}[Determinants and positive connections]\label{lem:det}
For the normalized fundamental matrix $Y$ of a positive connection,
\[
 d(\det Y)=(\tr\Omega)\det Y.
\]
Consequently a traceless positive connection has determinant-one transport
along every path.
\end{lemma}
\begin{proof}
The determinant is a finite polynomial in the entries of $Y$.
The usual adjugate identity and Leibniz rule therefore apply coefficientwise,
and $Y$ is invertible by Theorem~\ref{thm:fundamental}. Jacobi's formula
reduces to the displayed identity. If $\tr\Omega=0$, the determinant is
constant along the normalized solution and equals $1$ at the starting
point. All identities involve finite coefficientwise algebra or the already
justified support products.
\end{proof}

\begin{theorem}[Traceless normalized realization]\label{thm:SL}
Suppose an admissible representation in Theorem~\ref{thm:RH} takes values in
$\SL_r(\K)$. For an entrywise complex-linear period section, its unique
normalized realizer satisfies $\tr\Omega^\sigma=0$. Conversely every
traceless positive connection has monodromy in $\SL_r(\K)$.
The unique framed gauge between two traceless connections with the same
monodromy has determinant one. These statements retain all the support and
logarithmic-extension conclusions above.
\end{theorem}
\begin{proof}
Use the recursion \eqref{eq:RH-recursion}. Assume every lower coefficient
$\Omega_\delta$ is traceless. Then the partial connection
$\Omega^{<\gamma}$ is traceless, so its monodromy $N_d$ has determinant one
by Lemma~\ref{lem:det}. By the construction, $N_d$ and the target $M_d$
agree below $\gamma$. Both have constant term $I_r$.

In the determinant polynomial, a difference at weight $\gamma$ can
contribute to that weight only with all other factors of weight zero.
Therefore
\[
 \coeff{\gamma}(\det M_d-\det N_d)
   =\tr\bigl(m_{d,\gamma}-h_{d,\gamma}^{<\gamma}\bigr).
\]
The left side is zero. Hence the defect vector in \eqref{eq:defect} is
traceless at every puncture. Since the same linear section is used for each
matrix entry,
\[
 \tr\Omega_\gamma
 =\sigma\bigl((\tr b_{d,\gamma})_{d\in D}\bigr)=0.
\]
This proves the induction. The converse is Lemma~\ref{lem:det}.
For a gauge satisfying $dG=\Psi G-G\Omega$, Jacobi's formula gives
$d\det G=(\tr\Psi-\tr\Omega)\det G=0$.
Its value at $p$ is $1$, proving the final assertion.
\end{proof}

\section{Sharp global obstructions and distinguishing examples}
\label{sec:obstructions}

\subsection{Finite solvability does not imply global solvability}
\begin{example}[A countable special-linear obstruction]\label{ex:descending}
Let $D=\{1,2,3,\ldots\}$, let $\Gamma=\Q$, and set
\[
 N=\begin{pmatrix}0&1\\0&0\end{pmatrix},\qquad
 M_n=I_2+t^{1/n}N.
\]
Because the meridians form a free basis, this prescription defines a
representation $\rho:\pi_1(\C\setminus D,p)\to\SL_2(\K)$. Each generator
matrix has one positive exponent, is invertible, and is an entirely
legitimate Hahn matrix. Every finite initial collection of the prescribed
monodromies is realizable on the corresponding finitely punctured plane by
Corollary~\ref{cor:finite}. Nevertheless no positive globally Hahn-supported
holomorphic connection on $\C\setminus D$ has all the prescribed monodromies.
\end{example}
\begin{proof}
The determinant is $1$ because $N^2=0$ and $N$ is strictly upper triangular.
A finite set of the exponents $1/n$ is well ordered. The full union
\[
 E_\rho=\{1,1/2,1/3,\ldots\}
\]
is not well ordered: it is a strictly decreasing sequence and has no least
member. The nonexistence follows from Theorem~\ref{thm:RH}, even without any
requirement of simple poles. Proposition~\ref{prop:conjugation} shows that a
fixed change of matrix basis or an enlargement of the value group cannot
repair the obstruction.
\end{proof}

One-puncture realizers can even be written down:
\[
 \Omega_n=\frac{t^{1/n}}{2\pi i}N\frac{dz}{z-n},
 \qquad
 \exp\left(\int_{\ell_n}\Omega_n\right)=I_2+t^{1/n}N.
\]
The failure is not an inability to solve a differential equation at any one
puncture. It is the impossibility of placing the entire prescription in one
global support set. This is a genuine local-to-global obstruction in the
chosen coefficient category, not a claim of nonexistence of an ordinary
local system: the representation itself already defines such an algebraic
local system over $\K$.

\subsection{The union condition must not be strengthened to point-finiteness}
\begin{example}[Infinitely many punctures at one exponent]
\label{ex:same-exponent}
Let $D=\Z$, $U=\C\setminus\Z$, $\Gamma=\Q$, and use the same nilpotent
matrix $N$. Then
\begin{equation}\label{eq:cot}
 \Omega=\frac{t}{2i}\,N\cot(\pi z)\,dz
\end{equation}
has monodromy $I_2+tN$ at every positive based meridian.
The family $(M_n-I_2)_{n\in\Z}$ is \emph{not} strongly summable as a family
indexed by punctures, because exponent $1$ appears infinitely many times.
It is nevertheless globally realizable.
\end{example}
\begin{proof}
The ordinary form $\cot(\pi z)\,dz$ has residue $1/\pi$ at every integer.
Thus each meridian period of \eqref{eq:cot} is $tN$. Every product of two
coefficient matrices is zero, since both are multiples of $N$ and $N^2=0$.
The Dyson expansion therefore terminates after its length-one term, giving
$I_2+tN$. The connection itself has the singleton support $\{1\}$ and is
holomorphic on one common domain $U$. Nothing in the definition asks that
its periods over \emph{different} loops form a summable family.
\end{proof}

Together, Examples~\ref{ex:descending} and \ref{ex:same-exponent} identify
the precise strength of the main condition. Requiring only individual Hahn
matrices is too weak. Requiring a strongly summable family over punctures
is too strong. A well-ordered \emph{union} of supports is exactly right.

\subsection{Admissible supports need not be locally finite in an interval}
\begin{example}[Infinitely many exponents below one bound]
\label{ex:bounded-support}
On $D=\{1,2,\ldots\}$ prescribe
\[
 M_n=I_2+t^{2-1/n}N.
\]
The support $E=\{2-1/n:n\geq1\}$ is increasing of order type $\omega$, hence
well ordered, although it has infinitely many elements below $2$.
These data are realizable. More explicitly, for a fixed period section put
\[
 \alpha_n=\sigma(e_n),\qquad
 \Omega=N\sum_{n\geq1}t^{2-1/n}\alpha_n,
\]
where $e_n\in\C^D$ is $1$ at $n$ and $0$ elsewhere. The connection has support
in $E$, and its holonomy is exactly the prescribed family because $N^2=0$.
\end{example}

The displayed sum is a Hahn family, not an assertion that its ordinary
partial sums converge in the valuation topology. In $\Q$, the tail
valuations $2-1/n$ stay below $2$, so those partial sums do not even approach
the full family in that topology. This example also shows why a procedure
that promises to enumerate all coefficients below a fixed bound in finitely
many steps is not supplied by our theorem.

\subsection{Why logarithmic hypotheses matter for entire gauges}
\begin{example}[A nonlogarithmic connection with trivial monodromy]
On $\C^*$ consider the scalar connection
$\Omega=-t\,dz/z^2$ and base point $p=1$. Its normalized solution is
$Y=\exp(t(1/z-1))$, a positive globally supported holomorphic family on
$\C^*$. It is single-valued, so the monodromy is trivial. The unique framed
gauge from $\Omega$ to the zero connection is
$G=Y^{-1}=\exp(-t(1/z-1))$. Its coefficient at $t$ has a pole at $0$.
Therefore the entire-extension assertion of
Theorem~\ref{thm:entire-gauge} cannot be extended to connections with
arbitrary higher-order coefficient poles.
\end{example}

\section{A noncommuting calculation across two valuation ranks}
\label{sec:example}

\subsection{The targets and the first-order guess}
Take $U=\C\setminus\{0,1\}$ and base point $b=1/2$. Choose the usual
positive based meridians: a real tail to a small circle about the puncture,
a positive circuit, and the reversed tail. Define
\[
 X=\begin{pmatrix}0&1\\0&0\end{pmatrix},\quad
 Y=\begin{pmatrix}0&0\\1&0\end{pmatrix},\quad
 H=[X,Y]=\begin{pmatrix}1&0\\0&-1\end{pmatrix},
\]
so $X^2=Y^2=0$. Initially use two independent formal bookkeeping variables
$u,v$ of positive degree and prescribe
\begin{equation}\label{eq:two-targets}
 M_0=I_2+uX,\qquad M_1=I_2+vY.
\end{equation}
At the end set $u=t^\alpha$ and $v=t^\beta$ for positive, rationally
independent $\alpha,\beta\in\Gamma$. The choice $\alpha=1$, $\beta=\omega$
is the rank-two surcomplex example.

Put
\[
 \eta_0=\frac{dz}{2\pi i z},\qquad
 \eta_1=\frac{dz}{2\pi i(z-1)},\qquad
 C=\frac{\log 2}{2\pi i}.
\]
The obvious first-order guess is
\begin{equation}\label{eq:linear-guess}
 \Omega^{(1)}=uX\eta_0+vY\eta_1.
\end{equation}
Its first-order periods are correct. Its mixed-order monodromy is not.

\begin{proposition}[The first nonabelian correction]\label{prop:mixed}
For the connection \eqref{eq:linear-guess}, the coefficient of $uv$ in the
based monodromy is $CH$ at the meridian of $0$ and $-CH$ at the meridian
of $1$. Consequently the unique normalized realizer has expansion
\begin{equation}\label{eq:mixed-connection}
 \boxed{\quad
 \Omega=uX\eta_0+vY\eta_1
          -uv\,C H(\eta_0-\eta_1)+O_{\deg\geq3}.
 \quad}
\end{equation}
Its coefficients of $u^2$ and $v^2$ are zero. Here the remainder notation
refers only to total degree in $u,v$, not to a claimed valuation bound.
\end{proposition}

\begin{proof}
Along a chosen meridian let
$L_j(s)=\int_0^s\eta_j$ denote the accumulated integral from the base point.
The first-order transport is $I_2+uXL_0+vYL_1$. The mixed term in the
coefficient equation is therefore
\begin{equation}\label{eq:mixed-integral}
 XY\int_\ell L_1\eta_0+YX\int_\ell L_0\eta_1.
\end{equation}
For the meridian of $0$, $L_1$ is the single-valued ordinary function
\[
 L_1(z)=\frac{\log(1-z)-\log(1-b)}{2\pi i}
\]
on a neighborhood containing the tail and the disk about $0$ but not $1$.
It satisfies $L_1(b)=0$ and $L_1(0)=C$.
Integration of $d(L_0L_1)$ along the full loop gives
\[
 \int_\ell L_0\eta_1=-\int_\ell L_1\eta_0,
\]
because the product vanishes at both endpoints: $L_0$ ends at $1$ but
$L_1$ ends at $0$. The residue theorem gives
$\int_\ell L_1\eta_0=L_1(0)=C$; contributions of the outgoing and incoming
tails cancel. Thus \eqref{eq:mixed-integral} equals $C(XY-YX)=CH$.

For the meridian of $1$, use instead
\[
 L_0(z)=\frac{\log z-\log b}{2\pi i}.
\]
It is single-valued near that tail and disk, has $L_0(b)=0$, and
$L_0(1)=C$. Now
$\int_\ell L_0\eta_1=C$ and
$\int_\ell L_1\eta_0=-C$, so the mixed coefficient is $-CH$.
The prescribed matrices \eqref{eq:two-targets} have zero $uv$ coefficient.
The ordinary period section for $\{0,1\}$ therefore assigns the correcting
one-form $-CH\eta_0+CH\eta_1$ at that degree, proving
\eqref{eq:mixed-connection}. Pure second-order terms in the uncorrected
transport vanish because $X^2=Y^2=0$, so their defects and normalized
correction coefficients are zero.
\end{proof}

In residue form, the same calculation is
\begin{align}
 R_0&=\frac{u}{2\pi i}X
      -\frac{uv\log 2}{(2\pi i)^2}H+O_{\deg\geq3},\label{eq:R0}\\
 R_1&=\frac{v}{2\pi i}Y
      +\frac{uv\log 2}{(2\pi i)^2}H+O_{\deg\geq3}.\label{eq:R1}
\end{align}
The mixed corrections cancel in the residue at infinity, but are essential
at the two finite poles. All displayed coefficient matrices are traceless,
in agreement with Theorem~\ref{thm:SL}.

\subsection{The representation is genuinely nonabelian}
The exact group commutator of the prescribed matrices is
\begin{equation}\label{eq:commutator}
 M_0M_1M_0^{-1}M_1^{-1}
 =\begin{pmatrix}
  1+uv+u^2v^2&-u^2v\\
  uv^2&1-uv
 \end{pmatrix}.
\end{equation}
Its determinant is $1$, and its leading mixed term is $uvH\ne0$.
Thus the realization theorem here is not just an entrywise construction of
independent scalar equations. The extra commutator in
\eqref{eq:mixed-connection} is exactly the interaction that scalar period
prescription would miss.

\subsection{Specialization to a genuinely higher-rank Hahn group}
Set $\Gamma=\Q+\Q\omega$, $u=t$ and $v=t^\omega$. The monoid generated by
$1$ and $\omega$ is
\[
 \{a+b\omega:a,b\in\N\cup\{0\}\},
\]
ordered first by $b$ and then by $a$. It is well ordered. The correction in
\eqref{eq:mixed-connection} occurs at exponent $1+\omega$.
No finite number of increments by $1$ can reach even $\omega$, but the
well-founded recursion handles $\omega$ and $1+\omega$ directly. There is
no change in the proof or its support certificate.

The notation $O_{\deg\geq3}$ is especially important here. A term $t^3$
has total degree three and yet has smaller valuation than $t^\omega$.
Conversely a term $t^{3\omega}$ has much larger valuation. A finite
bivariate degree truncation is therefore not an approximation uniformly
valid below a single Hahn valuation threshold. The example checks a finite
coefficient identity, not a truncation of all coefficients below
$1+\omega$.

\section{Effectivity, base change, and a necessary discontinuity}
\label{sec:effectivity}

\subsection{What the coefficient recursion actually depends on}
\begin{proposition}[Finite exponent dependence]\label{prop:dependence}
Fix an admissible support $E$, a period section $\sigma$, a target exponent
$\gamma\in E^*\setminus\{0\}$, and the meridian paths. The coefficient
$\Omega^\sigma_\gamma$ is determined by finitely many exponent-indexed
vectors of target coefficients
\[
 (m_{d,\delta})_{d\in D},\qquad \delta\in E,
\]
and finitely many exponent patterns of ordinary products, iterated
integrals, and applications of $\sigma$. The set of punctures within one
such vector may still be infinite.
\end{proposition}
\begin{proof}
Every lower coefficient entering the defect at $\gamma$ has a weight that
occurs as a part in a decomposition of $\gamma$ into members of $E^*$.
There are finitely many such decompositions and parts by
Lemma~\ref{lem:finite-closure}. Recursively replacing a constructed
coefficient by its own defining defect only refines one of these finite
exponent-word decompositions. Equivalently, all expressions can be
organized by finite rooted trees whose leaves are $E$-letters summing to
$\gamma$. There are finitely many possible leaf words by Neumann's lemma,
and finitely many tree shapes and partitions for each fixed finite number
of leaves. Thus only finitely many exponent-level operations are involved.
Applying $\sigma$ may use an entire element of $\C^D$, so this argument
says nothing about finite dependence on puncture coordinates.
\end{proof}

A relative algorithm is now explicit. Represent the support well enough to
list the finitely many words relevant to a requested exponent. Compute the
lower-weight iterated-integral prediction, subtract it from the target
coefficient vector, and apply the chosen ordinary period section.
The theorem does not imply that these steps are computable for every
set-theoretically specified Hahn series, nor that a non-effective support
can be enumerated, nor that an infinite product vector can be supplied to a
finite program. For finite puncture sets the period section is explicit,
and the remaining tasks are ordinary symbolic or numerical integrations.

\begin{proposition}[Compatibility with enlargement of the workspace]
\label{prop:base-change}
Let $\iota:\Gamma\hookrightarrow\Gamma'$ be an order-preserving group
embedding, and use the same complex coefficient data, meridians, and period
section. For an admissible representation, its normalized realizer in the
larger Hahn workspace is obtained from the original realizer by replacing
$t^\gamma$ with $t^{\iota(\gamma)}$. The same compatibility holds for
normalized fundamental matrices and framed comparison gauges.
\end{proposition}
\begin{proof}
Order embeddings preserve Hahn summability and addition of exponents.
The finite coefficient formulas for products, iterated integrals, and
ordinary primitives are unchanged. The transfinite recursion on $E^*$
therefore maps to the same recursion on $\iota(E)^*$, with no extra
coefficients outside that monoid. Alternatively, the transported family
satisfies all defining equations and normalizations, so uniqueness in
Theorems~\ref{thm:fundamental}, \ref{thm:RH}, and \ref{thm:gauge} identifies
it with the larger-workspace construction.
\end{proof}

\subsection{Why an infinite linear period section cannot be continuous}
The choice-dependent period section in Lemma~\ref{lem:period-section} is
not merely missing an estimate that can automatically be supplied. A
simple analytic obstruction rules out the most natural continuity demand.
This is an auxiliary observation, not a claim of novelty for the general
functional-analytic phenomenon.

Give $\C^D$ its product topology. Give holomorphic one-forms on $U$ their
usual topology of uniform convergence of the coefficient functions on
compact subsets. Restrict this topology to $\LL(D)$.

\begin{proposition}[No continuous linear period right inverse]
\label{prop:discontinuous}
If $D$ is infinite, there is no continuous complex-linear right inverse
$\sigma:\C^D\to\LL(D)$ to $P$ with these topologies. In fact no such
continuous linear right inverse exists even if its values are allowed to
be arbitrary holomorphic one-forms on $U$ with the specified meridian
periods.
\end{proposition}
\begin{proof}
Choose a closed ordinary disk $K\subset U$ with nonempty interior. Write a
one-form on $U$ as $f(z)\,dz$, and let
$q_K(f\,dz)=\sup_{z\in K}|f(z)|$.
If $\sigma$ were continuous and linear, the continuous seminorm
$q_K\circ\sigma$ on the product $\C^D$ would be bounded by a constant times
a maximum of finitely many coordinate absolute values. Thus there would
be a finite set $F\subset D$ and $C_0>0$ such that
\[
 q_K(\sigma(v))\leq C_0\max_{d\in F}|v_d|\qquad(v\in\C^D).
\]
This standard product-topology bound also follows directly by taking a
basic finite-coordinate neighborhood on which the seminorm is less than
one and then scaling.

Choose $d\in D\setminus F$ and take the coordinate vector $e_d$.
The bound implies $q_K(\sigma(e_d))=0$. The coefficient function of
$\sigma(e_d)$ vanishes on a disk, so the ordinary identity theorem on the
connected domain $U$ makes it identically zero. Its period around $\ell_d$
is then zero, contradicting $(P\sigma(e_d))_d=1$.
\end{proof}

Thus the algebraic splitting used in the infinite-puncture normal form
cannot have product-topology continuity. This does not contradict the
existence theorem, whose coefficient category imposes no such continuity.
It also does not rule out nonlinear choices with other regularity
properties, or useful linear choices on smaller weighted sequence spaces.
Those are different problems. For finite $D$ the explicit section
\eqref{eq:finite-section} is continuous, and the obstruction disappears.

\section{Verification of the finite example}\label{sec:verification}

The file \texttt{verify\_examples.py} accompanying this article has two
independent components. SymPy checks the nilpotence identities, the
commutator $[X,Y]$, both target determinants, the exact matrix formula
\eqref{eq:commutator}, its determinant, and the trace of the correction.
These are eight exact algebraic checks.

A second component integrates the coefficient equations through bivariate
total degree two. It does not substitute small floating-point numbers for
$u$ and $v$. Instead it stores the matrices at degrees
\[
 (0,0),\ (1,0),\ (0,1),\ (2,0),\ (1,1),\ (0,2)
\]
and solves the finite triangular ODE system
\[
 \frac{dY_{a,b}}{ds}
    =\sum_{(c,d)\ne(0,0)}A_{c,d}(s)Y_{a-c,b-d}(s).
\]
For the meridian of $0$, the path goes from $1/2$ to $1/4$, follows the
circle of radius $1/4$ positively around $0$, and returns. For the meridian
of $1$, it goes from $1/2$ to $3/4$, follows the radius-$1/4$ circle about
$1$ from angle $\pi$ to $3\pi$, and returns.
The uncorrected mixed term is compared with $CH$ and $-CH$ respectively;
the corrected transport is compared with the target through degree two.

The recorded run used NumPy 2.3.5, SciPy 1.17.0, and SymPy 1.14.0. Its
maximum entrywise errors were as follows; these are observed numerical
errors, not rigorous error enclosures.
\begin{center}
\begin{tabular}{lrr}
\toprule
Check & Meridian at $0$ & Meridian at $1$\\
\midrule
Uncorrected mixed-term formula & $1.17\cdot10^{-13}$ & $1.17\cdot10^{-13}$\\
Corrected transport, degree $\leq2$ & $1.47\cdot10^{-13}$ & $1.47\cdot10^{-13}$\\
\bottomrule
\end{tabular}
\end{center}
The source and full-precision results are included in the archive. The
checks confirm the sign and normalization in the displayed finite example.
They do not test an arbitrary Hahn field, justify a transfinite recursion,
prove a support set well ordered, verify the theorems in Lean, or establish
that the results are unpublished. The mathematical proofs do not rely on
the numerical tests.

\section{Scope, unresolved extensions, and conclusion}\label{sec:limits}

\subsection{Exactly what has been established}
We have constructed an inverse to holonomy in a specified global
common-domain Hahn category, with a necessary-and-sufficient support
condition for an arbitrary closed discrete puncture set in the ordinary
plane. The main theorem supplies coefficientwise logarithmic realizers,
a support certificate, and a unique representative for a chosen linear
normalization. The gauge theorem identifies precisely which positive
connections carry the same framed monodromy. Within the logarithmic class
the comparison gauge extends across every finite puncture. The determinant
refinement and the noncommuting example show that the result includes
nonabelian special-linear data, not just scalar additive interpolation.

The construction is genuinely sensitive to global support. It explains both
why infinitely many individually valid matrices can be impossible to realize
and why infinitely many punctures at one and the same exponent are
perfectly allowable. Enlarging the surreal workspace does not change either
conclusion.

\subsection{What is deliberately not asserted}
Positivity is essential to the present proofs. We have not treated arbitrary
nontrivial exponent-zero monodromy, arbitrary ordinary residues with
resonances, arbitrary algebraic groups, or a prescribed condition at
infinity for an infinite puncture set. We have not proved a general
Riemann--Hilbert correspondence on every surcomplex domain. The ordinary
base, ordinary universal cover, and ordinary path category are part of the
statement, not provisional notation for a fine-topological construction.

We have also not imposed norm convergence in the Hahn parameter, a growth
condition on coefficient functions, or a topological completeness assumption
on $\Gamma$. Adding such conditions can change the theorem. In particular,
Proposition~\ref{prop:discontinuous} prevents one natural continuous linear
normalization for unrestricted infinite residue data. The freedom to use
arbitrary ordinary entire correction terms at each exponent is important to
the infinite-puncture existence proof.

No theorem here relies on an unrefereed companion claim about Hahn
Noetherianity, a Nullstellensatz, or global sheaf cohomology. The imported
inputs are the classical support lemma, ordinary primitives on simply
connected Riemann surfaces, elementary topology of punctured planar
domains, and ordinary Mittag--Leffler. These dependencies make the proof
reasonably isolated for independent checking.

\subsection{Three concrete next problems}
\paragraph{Deformations of a nontrivial ordinary connection.}
Fix an ordinary connection $\Omega_0$ and a representation reducing to its
monodromy rather than to the identity. The first correction operator is
then a twisted period map with values in the adjoint local system, not the
scalar period map $P$. A precise extension should specify when that map has
a right inverse within a chosen pole class and should track any local
resonances. Our triangular proof does not by itself establish such a
right inverse.

\paragraph{Controlled growth and effective infinite data.}
For a prescribed weight or growth condition on residues at $D$, identify
spaces on which a useful continuous or computable period section exists.
The exact question is to characterize the residue sequence spaces admitted
by the chosen coefficient growth class, and then prove that nonlinear
monodromy corrections stay inside them. Arbitrary algebraic splittings
need not preserve any such structure.

\paragraph{Higher-dimensional bases.}
On a complex manifold of dimension at least two, a proposed correction
must satisfy the nonlinear flatness equation as well as the period
conditions. The proof here uses the automatic flatness of holomorphic
one-forms on a curve. A higher-dimensional analogue must account for
relations in the fundamental group and for deformation obstructions;
it cannot be inferred by changing the symbol $z$ to a tuple.

\subsection{Conclusion}
The support-sensitive problem has a complete answer within the stated
category: global well-ordering of the monodromy support is the exact
criterion, and a fixed period section produces an explicit recursive normal
form. The order of the value group may have several ranks, the singularity
set may be infinite, and the monodromy may be noncommutative. None of these
requires a false topological limiting argument. The analytic work is done
coefficientwise on the ordinary base; the Hahn support calculus is exactly
what makes those ordinary constructions coexist in one surcomplex
workspace.

\appendix
\section{A proof-dependency and reproducibility audit}\label{app:audit}

\subsection{Dependency graph}
The classical Neumann lemma implies the support control for inverses and
Dyson expansions. Ordinary primitives then give the fundamental-matrix
theorem and the triangular variation lemma. Ordinary Mittag--Leffler gives
the period section. Those two branches meet in the realization recursion.
The gauge classification uses only the fundamental matrices; its entire
extension uses only single-valuedness and the fact that a simple pole with
zero residue is removable. The traceless refinement adds the finite
polynomial determinant identity. The counterexamples use the necessity
half of the main theorem and explicit support orderings.

In particular, none of the implications depends on the outcome of the
numerical verification script. The script's role is confined to auditing
formulas \eqref{eq:mixed-connection} and \eqref{eq:commutator}.

\subsection{Literature and repository audit boundaries}
The repository catalogue and the README descriptions cited below were
retrieved at the pinned commit in Section~\ref{sec:intro}. The mathematical
source papers consulted for the Hahn background include Poonen and
M\"uller--Strohmaier; the latter's definition and positive-support inversion
were also inspected in the typeset paper. Classical differential-equation
references were used to locate the established regular-singular and
Riemann--Hilbert context. Bibliographic access to a book or its contents is
not described as an exhaustive inspection of every theorem in it.

Searches included combinations of ``Hahn'', ``monodromy'',
``Riemann--Hilbert'', and ``surcomplex'', together with searches for the
specific primary sources in the bibliography. The search was targeted,
not a proof of absence from the literature. The repository's incomplete
code-search response was not used to conclude that a term or theorem is
absent. The date of this assessment is September 21, 2026.

\subsection{Rebuilding and rerunning}
The LaTeX source is standalone and uses an internal bibliography. A standard
TeX installation can build it with
\begin{quote}\ttfamily\small
latexmk -pdf -interaction=nonstopmode -halt-on-error article.tex
\end{quote}
The verifier requires Python with NumPy, SciPy, and SymPy. Run
\begin{quote}\ttfamily\small
python verify\_examples.py
\end{quote}
It prints the results and writes \texttt{verification\_results.json} next to
the script. Its numerical assertions use an absolute tolerance of
$2\cdot10^{-9}$; the actual recorded errors are reported in
Section~\ref{sec:verification}. Floating-point library or platform changes
can alter those last digits without affecting the exact checks.

\begin{thebibliography}{99}
\small\raggedright
\addcontentsline{toc}{section}{References}

\bibitem{neumann}
B.~H. Neumann,
\emph{On ordered division rings},
Transactions of the American Mathematical Society \textbf{66} (1949),
202--252.
\href{https://www.jstor.org/stable/1990552}{Primary journal record}.

\bibitem{poonen}
B. Poonen,
\emph{Maximally complete fields},
L'Enseignement Math\'ematique (2) \textbf{39} (1993), 87--106.
DOI: \href{https://doi.org/10.5169/seals-60414}{10.5169/seals-60414}.
\href{https://math.stanford.edu/~conrad/Perfseminar/refs/poonencomplete.pdf}{Accessible copy of the paper}.

\bibitem{muller}
J. M\"uller and A. Strohmaier,
\emph{The theory of Hahn-meromorphic functions, a holomorphic Fredholm
theorem, and its applications},
Analysis \& PDE \textbf{7} (2014), no.~3, 745--770.
DOI: \href{https://doi.org/10.2140/apde.2014.7.745}{10.2140/apde.2014.7.745}.
\href{https://arxiv.org/abs/1205.0236}{arXiv:1205.0236}.

\bibitem{bournez}
O. Bournez and Q. Guilmant,
\emph{Surreal fields stable under exponential and logarithmic functions},
preprint (2022), \href{https://arxiv.org/abs/2201.08199}{arXiv:2201.08199}.
Cited for the normal-form and set-sized-surreal-field background, not for
the monodromy results of this article.

\bibitem{chen}
K.-T. Chen,
\emph{Iterated path integrals},
Bulletin of the American Mathematical Society \textbf{83} (1977),
831--879.
DOI: \href{https://doi.org/10.1090/S0002-9904-1977-14320-6}{10.1090/S0002-9904-1977-14320-6}.

\bibitem{deligne}
P. Deligne,
\emph{\'Equations diff\'erentielles \`a points singuliers r\'eguliers},
Lecture Notes in Mathematics, vol.~163, Springer, 1970.
DOI: \href{https://doi.org/10.1007/BFb0061194}{10.1007/BFb0061194}.

\bibitem{ilyashenko}
Y. Ilyashenko and S. Yakovenko,
\emph{Lectures on Analytic Differential Equations},
Graduate Studies in Mathematics, vol.~86, American Mathematical Society,
2008.
DOI: \href{https://doi.org/10.1090/gsm/086}{10.1090/gsm/086}.

\bibitem{repo-catalogue}
\texttt{VladimirReshetnikov/Surreal},
\emph{The surreal and surcomplex reports}, documentation catalogue,
\texttt{docs/README.md}, at commit \commit.
\href{https://github.com/VladimirReshetnikov/Surreal/blob/39f2be6667ade51bca2b45daa47e289d69c09764/docs/README.md}{Pinned source}.
Accessed September 21, 2026. The following four repository references use
the same commit and access date.

\bibitem{repo-analysis}
\texttt{VladimirReshetnikov/Surreal},
\emph{Surcomplex analysis: holomorphic functions on $K=\No[i]$},
\texttt{docs/surcomplex/analysis/README.md}.
\href{https://github.com/VladimirReshetnikov/Surreal/blob/39f2be6667ade51bca2b45daa47e289d69c09764/docs/surcomplex/analysis/README.md}{Pinned source}.
AI-assisted, unrefereed research documentation.

\bibitem{repo-geometry}
\texttt{VladimirReshetnikov/Surreal},
\emph{Infinitesimal analytic geometry over the surcomplex numbers:
radius-free germs, a Nullstellensatz, and finite maps},
\texttt{docs/surcomplex/analytic-geometry/README.md}.
\href{https://github.com/VladimirReshetnikov/Surreal/blob/39f2be6667ade51bca2b45daa47e289d69c09764/docs/surcomplex/analytic-geometry/README.md}{Pinned source}.
AI-assisted, unrefereed research documentation.

\bibitem{repo-global}
\texttt{VladimirReshetnikov/Surreal},
\emph{Global support obstructions: moving divisors, Mittag--Leffler, and a
Picard dichotomy},
\texttt{docs/surcomplex/global-divisors/README.md}.
\href{https://github.com/VladimirReshetnikov/Surreal/blob/39f2be6667ade51bca2b45daa47e289d69c09764/docs/surcomplex/global-divisors/README.md}{Pinned source}.
AI-assisted, unrefereed research documentation.

\bibitem{repo-contours}
\texttt{VladimirReshetnikov/Surreal},
\emph{Contours and Stokes: Jordan separation, Cauchy theory, and Stokes
calculus on the surcomplex plane},
\texttt{docs/surcomplex/contours-and-stokes/README.md}.
\href{https://github.com/VladimirReshetnikov/Surreal/blob/39f2be6667ade51bca2b45daa47e289d69c09764/docs/surcomplex/contours-and-stokes/README.md}{Pinned source}.
AI-assisted, unrefereed research documentation.

\end{thebibliography}
\end{document}
